\subsection{内自同构}

设 \(G\) 为一个群。\(\operatorname{Aut}(G)\) （群 \(G\) 的自同构所成的集合）是 \(\mathfrak{S}_{\mathbb{G}}\) 的一个子群（§ 4，第 1 号，例 2）。

命题 3. 设 G 为一个群。对 G 的每个元素 x，映射 \(\operatorname{Int}(x): y \mapsto xyx^{-1}\) 是 G 到自身的一个自同构。映射 \(\operatorname{Int}: x \mapsto \operatorname{Int}(x)\) 是 G 到 \(\operatorname{Aut}(G)\) 的群同态，其核为 G 的中心，其像为 \(\operatorname{Aut}(G)\) 的一个正规子群。

若 x、y 和 z 是 G 的元素，则 \((xyx^{-1})(xzx^{-1}) = xyzx^{-1}\) ，因此 \(\operatorname{Int}(x)\) 是 G 的一个自同态。对于 G 的元素 x 和 y，

\[
\operatorname{Int} (x) \circ \operatorname{Int} (y) = \operatorname{Int} (x y):
\]

对所有 \(z \in G\) ，有 \(x(yzy^{-1})x^{-1} = (xy)z(xy)^{-1}\) 。另一方面， \(\operatorname{Int}(e)\) 是 G 的恒等映射。因此映射 Int 是从 G 到群 G 的自同态幺半群 \(\operatorname{End}(G)\) 的幺半群同态。由于 G 的元素都可逆，映射 Int 取值于 \(\operatorname{Aut}(G)\) ，即 \(\operatorname{End}(G)\) 的可逆元素所成的集合（§ 2，第 3 号）。现在， \(xyx^{-1} = y\) 当且仅当 x 与 y 可交换，因此 \(\operatorname{Int}(x)\) 是 G 的恒等映射当且仅当 x 是中心元素。最后，设 \(\alpha\) 为 G 的一个自同构，并设 \(x \in G\) ；于是

\[
\operatorname{Int} (\alpha (x)) = \alpha \circ \operatorname{Int} (x) \circ \alpha^ {- 1}.\tag{2}
\]

对于 \(y \in G\) ,

\[
\alpha (x). y. \alpha (x) ^ {- 1} = \alpha (x). \alpha (\alpha^ {- 1} (y)). \alpha (x) ^ {- 1} = \alpha (x. \alpha^ {- 1} (y). x ^ {- 1}).
\]

因此 \(\alpha\) . Int(G) . \(\alpha^{-1} \subset \operatorname{Int}(G)\) 。

定义 4. 设 G 为一个群，\(x \in \mathbf{G}\) 。自同构 \(y \mapsto xyx^{-1}\) 称为由 \(x\) 定义的 G 的内自同构，记作 Int \(x\) 。

对于 \(x, y \in G\) ，我们也记 \(x^y = y^{-1}xy = (\operatorname{Int} y^{-1})(x)\) 。

G 的一个子群是正规的，当且仅当它在 G 的所有内自同构下稳定（§ 4，第 4 号，定义 5）。G 的一个子群若在 G 的所有自同构下稳定，则称为特征子群。群 G 的中心是一个特征子群（公式 (2)）。

群 G 的中心未必在 G 的所有自同态下稳定（习题 22）。特别地，带算子群的中心未必是稳定子群。

命题 4. 设 G 为一个群，H 为 G 的特征（相应地，正规）子群，K 为 H 的特征子群。则 K 是 G 的特征（相应地，正规）子群。

G 的自同构（相应地，内自同构）在 H 上的限制是 H 的自同构，从而保持 K 不变。

设 G 为一个群，A ⊂ G，且 \(b \in G\) 。若 \(bAb^{-1} = A\) ，则称 b 正规化 A；若对所有 \(a \in A\) 都有 \(bab^{-1} = a\) ，则称 b 中心化 A。设 A 和 B 是 G 的子集；若 B 的每个元素都正规化（相应地，中心化）A，则称 B 正规化（相应地，中心化）A。

G 中正规化（相应地，中心化）A 的 \(g \in G\) 所成的集合称为 A 的正规化子（相应地，中心化子）（参见 §1，第 5 号，定义 9）；它常记为 \(N_{G}(A)\) 或简记为 \(N(A)\) （相应地，\(C_{G}(A)\) 或 \(C(A)\) ）。它是 G 的一个子群。当 A 是 G 的子群时，\(N_{G}(A)\) 可刻画为 G 中包含 A 且使 A 在其中为正规子群的最大子群。

注记。(1) 当 G 通过内自同构作用于自身时，A 的正规化子（相应地，中心化子）就是 A 的严格稳定化子（相应地，固定子）。特别地，中心化子是正规化子的正规子群。

\begin{enumerate}
\def\labelenumi{(\arabic{enumi})}
\setcounter{enumi}{1}
\tightlist
\item
  \(b \in \mathbf{G}\) 中使 \(b\mathbf{A}b^{-1} \subset \mathbf{A}\) 的元素所成的集合是 \(\mathbf{G}\) 的一个子幺半群，即使 \(\mathbf{A}\) 是 \(\mathbf{G}\) 的子群，它也不一定是 \(\mathbf{G}\) 的子群（习题 27）。
\end{enumerate}

